% GENERATED FILE. Edit canonical lessons, then run npm run generate:books.
% canonical-source-hash: fnv1a64:c90a4b58185923b2
% canonical-lessons: SA-C318-sena, SA-C318-yodhah, SA-C318-yuddham, SA-C318-prayati, SA-C318-yudhyate

\chapter{The Army Marches}
\label{ch:sa-the-army-marches}
\hlchaptermodality{318}

\begin{chapteropening}
\textbf{By the end of this chapter:} \emph{I can tell how an army set out and fought.}

Complete the last lesson of chapter 318: I can tell how an army set out and fought.
\end{chapteropening}

\section[\texorpdfstring{\sk{सेना} (senā)}{सेना (senā)}]{\sk{सेना} (senā) — an army}

\label{lesson:SA-C318-sena}



\begin{quote}
Before the new one: say the Sanskrit for a peace treaty, then the Sanskrit for a vassal.
\end{quote}

\subsection*{You'll want to know: \sk{सेना}}
\textbf{\sk{सेना}} — \emph{senā} — \textquotedblleft{}an army\textquotedblright{}.

\begin{quote}
\sk{सेना} \sk{नगरात्} \sk{निर्गच्छति।}
\end{quote}

\begin{quote}
\emph{senā nagarāt nirgacchati.}
\end{quote}

\begin{quote}
The army marches out of the city.
\end{quote}

A feminine noun.

\begin{cousinweb}
\textbf{\sk{सेना}} (\emph{senā}) comes from the root √si, \textquotedblleft{}to bind\textquotedblright{}, by MW's account; its oldest sense is a missile or spear, and its usual sense an army, an armed force (MW).
\end{cousinweb}

\begin{grammarlens}[title={What you've built}]
One more word for telling what happened.
\end{grammarlens}

\subsection*{Your turn}
\begin{itemize}
\raggedright
  \item \emph{Say it:} \emph{senā}
  \item \emph{Say it:} \emph{senā}, once more
  \item \emph{Say it:} \emph{senā nagarāt nirgacchati.} — The army marches out of the city.
  \item \emph{Recall:} say the Sanskrit for a peace treaty, then the Sanskrit for a vassal, then say \emph{senā} again
\end{itemize}

\begin{tcolorbox}[breakable,colback=teal!4,colframe=teal!35!black,title={Before you move on}]
What does \sk{सेना} mean? (\textquotedblleft{}An army\textquotedblright{}.) Say it once more.
\end{tcolorbox}
\section[\texorpdfstring{\sk{योधः} (yodhaḥ)}{योधः (yodhaḥ)}]{\sk{योधः} (yodhaḥ) — a warrior}

\label{lesson:SA-C318-yodhah}



\begin{quote}
Before the new one: say the Sanskrit for a vassal, then the Sanskrit for an army.
\end{quote}

\subsection*{You'll want to know: \sk{योधः}}
\textbf{\sk{योधः}} — \emph{yodhaḥ} — \textquotedblleft{}a warrior\textquotedblright{}.

\begin{quote}
\sk{योधः} \sk{हस्ते} \sk{खड्गं} \sk{धारयति।}
\end{quote}

\begin{quote}
\emph{yodhaḥ haste khaḍgaṃ dhārayati.}
\end{quote}

\begin{quote}
The warrior holds a sword in his hand.
\end{quote}

A masculine noun.

\begin{cousinweb}
\textbf{\sk{योधः}} (\emph{yodhaḥ}) comes from the root √yudh, \textquotedblleft{}to fight\textquotedblright{}; MW gives \emph{yodha} as a fighter, warrior, soldier (MW).
\end{cousinweb}

\begin{grammarlens}[title={What you've built}]
One more word for telling what happened.
\end{grammarlens}

\subsection*{Your turn}
\begin{itemize}
\raggedright
  \item \emph{Say it:} \emph{yodhaḥ}
  \item \emph{Say it:} \emph{yodhaḥ}, once more
  \item \emph{Say it:} \emph{yodhaḥ haste khaḍgaṃ dhārayati.} — The warrior holds a sword in his hand.
  \item \emph{Recall:} say the Sanskrit for a vassal, then the Sanskrit for an army, then say \emph{yodhaḥ} again
\end{itemize}

\begin{tcolorbox}[breakable,colback=teal!4,colframe=teal!35!black,title={Before you move on}]
What does \sk{योधः} mean? (\textquotedblleft{}A warrior\textquotedblright{}.) Say it once more.
\end{tcolorbox}
\section[\texorpdfstring{\sk{युद्धम्} (yuddham)}{युद्धम् (yuddham)}]{\sk{युद्धम्} (yuddham) — a war; a battle}

\label{lesson:SA-C318-yuddham}



\begin{quote}
Before the new one: say the Sanskrit for an army, then the Sanskrit for a warrior.
\end{quote}

\subsection*{You'll want to know: \sk{युद्धम्}}
\textbf{\sk{युद्धम्}} — \emph{yuddham} — \textquotedblleft{}a war; a battle\textquotedblright{}.

\begin{quote}
\sk{सीमायां} \sk{युद्धं} \sk{भवति।}
\end{quote}

\begin{quote}
\emph{sīmāyāṃ yuddhaṃ bhavati.}
\end{quote}

\begin{quote}
There is a war on the border.
\end{quote}

A neuter noun.

\begin{cousinweb}
\textbf{\sk{युद्धम्}} (\emph{yuddham}) comes from the root √yudh, \textquotedblleft{}to fight\textquotedblright{}; MW gives \emph{yuddha} as fought and, as a neuter noun, a battle, fight, war (MW).
\end{cousinweb}

\begin{grammarlens}[title={What you've built}]
One more word for telling what happened.
\end{grammarlens}

\subsection*{Your turn}
\begin{itemize}
\raggedright
  \item \emph{Say it:} \emph{yuddham}
  \item \emph{Say it:} \emph{yuddham}, once more
  \item \emph{Say it:} \emph{sīmāyāṃ yuddhaṃ bhavati.} — There is a war on the border.
  \item \emph{Recall:} say the Sanskrit for an army, then the Sanskrit for a warrior, then say \emph{yuddham} again
\end{itemize}

\begin{tcolorbox}[breakable,colback=teal!4,colframe=teal!35!black,title={Before you move on}]
What does \sk{युद्धम्} mean? (\textquotedblleft{}A war; a battle\textquotedblright{}.) Say it once more.
\end{tcolorbox}
\section[\texorpdfstring{\sk{प्रयाति} (prayāti)}{प्रयाति (prayāti)}]{\sk{प्रयाति} (prayāti) — sets out}

\label{lesson:SA-C318-prayati}



\begin{quote}
Before the new one: say the Sanskrit for a warrior, then the Sanskrit for a war.
\end{quote}

\subsection*{You'll want to know: \sk{प्रयाति}}
\textbf{\sk{प्रयाति}} — \emph{prayāti} — \textquotedblleft{}sets out\textquotedblright{}.

\begin{quote}
\sk{प्रातः} \sk{सेना} \sk{युद्धाय} \sk{प्रयाति।}
\end{quote}

\begin{quote}
\emph{prātaḥ senā yuddhāya prayāti.}
\end{quote}

\begin{quote}
In the morning the army sets out for battle.
\end{quote}

Past (it …): \sk{सेना} \sk{प्रातः} \sk{प्रयातवती}, \emph{senā prātaḥ prayātavatī}.

\begin{cousinweb}
\textbf{\sk{प्रयाति}} (\emph{prayāti}) is built from \emph{pra}, \textquotedblleft{}forth\textquotedblright{}, and the root √yā, \textquotedblleft{}to go\textquotedblright{}; MW gives \emph{pra-yā} as to go forth, set out, advance (MW).
\end{cousinweb}

\begin{grammarlens}[title={What you've built}]
One more word for telling what happened.
\end{grammarlens}

\subsection*{Your turn}
\begin{itemize}
\raggedright
  \item \emph{Say it:} \emph{prayāti}
  \item \emph{Say it:} \emph{prayāti}, once more
  \item \emph{Say it:} \emph{prātaḥ senā yuddhāya prayāti.} — In the morning the army sets out for battle.
  \item \emph{Recall:} say the Sanskrit for a warrior, then the Sanskrit for a war, then say \emph{prayāti} again
\end{itemize}

\begin{tcolorbox}[breakable,colback=teal!4,colframe=teal!35!black,title={Before you move on}]
What does \sk{प्रयाति} mean? (\textquotedblleft{}Sets out\textquotedblright{}.) Say it once more.
\end{tcolorbox}
\section[\texorpdfstring{\sk{युध्यते} (yudhyate)}{युध्यते (yudhyate)}]{\sk{युध्यते} (yudhyate) — fights}

\label{lesson:SA-C318-yudhyate}



\begin{quote}
Before the new one: say the Sanskrit for a war, then the Sanskrit for sets out.
\end{quote}

\subsection*{You'll want to know: \sk{युध्यते}}
\textbf{\sk{युध्यते}} — \emph{yudhyate} — \textquotedblleft{}fights\textquotedblright{}.

\begin{quote}
\sk{योधाः} \sk{शौर्येण} \sk{युध्यन्ते।}
\end{quote}

\begin{quote}
\emph{yodhāḥ śauryeṇa yudhyante.}
\end{quote}

\begin{quote}
The warriors fight with valour.
\end{quote}

Past (he …): \sk{सः} \sk{शौर्येण} \sk{युद्धवान्}, \emph{saḥ śauryeṇa yuddhavān}.

\begin{cousinweb}
\textbf{\sk{युध्यते}} (\emph{yudhyate}) comes from the root √yudh, \textquotedblleft{}to fight\textquotedblright{}, which also gives \emph{yuddha}, \textquotedblleft{}a war\textquotedblright{} (MW).
\end{cousinweb}

\begin{grammarlens}[title={What you've built}]
One more word for telling what happened.
\end{grammarlens}

\subsection*{Your turn}
\begin{itemize}
\raggedright
  \item \emph{Say it:} \emph{yudhyate}
  \item \emph{Say it:} \emph{yudhyate}, once more
  \item \emph{Say it:} \emph{yodhāḥ śauryeṇa yudhyante.} — The warriors fight with valour.
  \item \emph{Recall:} say the Sanskrit for a war, then the Sanskrit for sets out, then say \emph{yudhyate} again
\end{itemize}

\begin{tcolorbox}[breakable,colback=teal!4,colframe=teal!35!black,title={Before you move on}]
What does \sk{युध्यते} mean? (\textquotedblleft{}Fights\textquotedblright{}.) Say it once more.
\end{tcolorbox}
